\begin{code}[hide]
module vqupit.sections.permutations where

open import Data.Nat using (ℕ ; suc)
open import Data.Product using (_×_ ; _,_ ; proj₁ ; proj₂)
open import Data.Unit using (⊤)
open import Level using (0ℓ)
open import Relation.Binary using (Rel)

open import Notations using (₁₊ ; ₂₊ ; ₃₊)
open import Word.Base using (Word ; WRel ; [_]ʷ ; ε ; _•_ ; wmap)
open import Presentation.Definitions using (_IsPresentationOf_)
open import Examples.Groups.Symmetric.Semantics using (Permutation′-group)
import Examples.Groups.Symmetric.Presentation as SymPres

postulate
  proof-elided : ∀ {ℓ} {A : Set ℓ} → A

private variable
  n : ℕ
\end{code}
\section{An Introduction to the Framework}\label{sec:permutations}

The formalisation is built on a library for presentations, normal forms
and completeness of circuit relation
sets~\cite{bian2026framework,qupit-code}.  This section introduces it on
its smallest complete case study: circuits built from a single two-wire
gate, the swap $\sigma$, which present the symmetric groups.  The
walkthrough exhibits the pipeline --- gates,
relations, cosets, tower, semantics, uniqueness, completeness,
presentation --- that \S\ref{sec:symplectic} instantiates again for the
symplectic factor, concretely enough that the latter can be read as the
same pipeline with richer gates and a richer action.  The section closes
with the product constructions (\S\ref{sec:constructions}) by which
\S\ref{sec:composing} glues such presentations together.

\subsection{Gates, generators, and circuits}\label{sec:sn-gates}

A gate set is a family $\aid{Gate} : \mathbb{N} \to \mathsf{Set}$ giving,
for each arity, the gates on that many adjacent wires; for permutations
there is one gate, on two wires.  The generic circuit layer of the
library turns any such family into wire-indexed
\emph{generators}: a gate placed at the bottom of the wires, or a
generator shifted up by one wire.

\begin{code}
data Gate : ℕ → Set where
  σ-gate : Gate 2
data Gen : ℕ → Set where
  gate₁ : Gate 1 → Gen (₁₊ n)     -- a 1-ary gate on the bottom wire
  gate₂ : Gate 2 → Gen (₂₊ n)     -- a 2-ary gate on the bottom two wires
  _↥   : Gen n → Gen (₁₊ n)       -- the same generator, one wire higher
\end{code}
\begin{code}[hide]
Circuit : ℕ → Set
Circuit n = Word (Gen n)

CRel : ℕ → Set₁
CRel n = Rel (Circuit n) 0ℓ

_↑ : Circuit n → Circuit (₁₊ n)
_↑ = wmap _↥

pattern σ-gen = gate₂ σ-gate

σ : Circuit (₂₊ n)
σ = [ σ-gen ]ʷ

infix 4 _SRel,_===_ _VRel,_===_ _≈_
infixr 10 σ•_

private variable
  w v : Circuit n

postulate
  _≈_ : Circuit n → Circuit n → Set
\end{code}

Here $\aid{₁₊}\;n$ and $\aid{₂₊}\;n$ are the library's patterns for
$1 + n$ and $2 + n$, and $n$ is an implicit argument, so
$\aid{gate₂}\;\aid{σ-gate}$ is a generator at any width from two upward.
An $n$-wire \emph{circuit} is a word of generators, $\aid{Circuit}\;n =
\aid{Word}\;(\aid{Gen}\;n)$: a single generator
$[\,g\,]$\aid{ʷ}, the empty circuit \aid{ε} --- $n$ bare wires, the
identity circuit at that width --- or a sequential composition
$c \aidop{•} c'$.  We write $\sigma$ for the one-letter circuit
$[\,\aid{gate₂}\;\aid{σ-gate}\;]$\aid{ʷ} and $c\;{\uparrow}$ for $c$ shifted
up one wire ($\aid{{\_}↑} = \aid{wmap}\;\aid{{\_}↥}$).  Read
diagrammatically, with wires drawn bottom-up and a word drawn left to
right as written, $\sigma$ is a crossing of the bottom two wires, $\sigma\;{\uparrow}$
the same crossing one wire higher, and $\sigma \aidop{•} \sigma\;{\uparrow}
\aidop{•} \sigma$ the three-wire braid on the left of the
\aid{yang-baxter} equation drawn beside its axiom below.

\subsection{Relations: two axioms, structure for free}\label{sec:sn-relations}

The group-specific relations are an inductive family indexed by the wire
count.  Each axiom is stated once, with its gates at the bottom of the
wires and the width left open (the implicit $n$), so that one constructor
covers every width at which the axiom fits; the structural rule
\aid{cong↑} below moves it to any height.  (Agda names may be mixfix:
each underscore in \aid{{\_}SRel,{\_}==={\_}} marks an argument
position, so the family is applied as $n\;\aid{SRel,}\;w \aidrel{===}
v$, read ``$w = v$ is an axiom on $n$ wires''.)

\begin{center}
\begin{minipage}[c]{0.58\linewidth}
\begingroup\setlength{\mathindent}{0pt}
\begin{code}
data _SRel,_===_ : (n : ℕ) → WRel (Gen n) where
  order       : (₂₊ n) SRel, σ • σ === ε
  yang-baxter : (₃₊ n) SRel, σ • σ ↑ • σ === σ ↑ • σ • σ ↑
\end{code}
\endgroup
\end{minipage}\hspace{2.5em}%
\begin{minipage}[c]{0.28\linewidth}\centering
  \ufig[0.35]{sn-order}\\[4pt]
  \ufig[0.35]{sn-yb}
\end{minipage}
\end{center}

\noindent These two axioms are two thirds of the Coxeter presentation of
$S_n$: $\sigma^2 = \varepsilon$ and the braid (Yang--Baxter) relation.
The third Coxeter relation, that swaps on disjoint wires commute, is not
specific to permutations: it is one of the relations every circuit
theory needs, whatever its gates.  By a
\emph{circuit theory} we mean a gate set with gate-specific relations,
presenting for each width $n$ a monoid of $n$-wire circuits under
sequential composition.  Such circuits are the fixed-width slices of a
symmetric monoidal category (a PROP), and the monoidal structure leaves
exactly two traces at fixed width: the shift $c \mapsto c\;{\uparrow}$ is
``tensoring with an identity wire'', so equations must be preserved by
it, and the interchange law says that gates on disjoint wires commute.
(The library also admits $0$-ary gates --- global scalars, available at
every width --- with one more structural rule, that a scalar is
identified with its own shift, \aid{ω↑=ω}; that it then commutes with
every generator is a theorem, once scalars are assumed to commute with
one another, which is vacuous for a single scalar.)  These are the structural rules every
circuit theory shares --- what a
pen-and-paper presentation leaves implicit in ``a congruence compatible
with composition and tensor product'' --- and the circuit layer adds them
generically: \aid{Lift-Relation} extends any
gate-specific family to the full relation set \aid{{\_}VRel,{\_}==={\_}}:

\begin{code}
data _VRel,_===_ : (n : ℕ) → CRel n where
  srel  : n SRel, w === v → n VRel, w === v
  cong↑ : n VRel, w === v → (₁₊ n) VRel, w ↑ === v ↑
  comm₁ : (h : Gate 1) (g : Gen n) → (₁₊ n) VRel, [ g ↥ ]ʷ • [ gate₁ h ]ʷ === [ gate₁ h ]ʷ • [ g ↥ ]ʷ
  comm₂ : (h : Gate 2) (g : Gen n) → (₂₊ n) VRel, [ g ↥ ↥ ]ʷ • [ gate₂ h ]ʷ === [ gate₂ h ]ʷ • [ g ↥ ↥ ]ʷ
\end{code}

From here on,
$\approx$ denotes the monoid congruence generated by
$n\;\aid{VRel,{\_}==={\_}}$: the least relation
containing the axioms and closed under reflexivity, symmetry, transitivity,
congruence of \aidop{•}, associativity, and the unit laws; deciding it is
the word problem.  Because $\sigma$ is its own inverse and shifts of
invertible generators are invertible, the presented monoid is a group; we
witness this by a \aid{Grouplike} proof (every generator has a left
inverse), by induction on generators.

\subsection{Cosets and the staircase normal form}\label{sec:sn-cosets}

The engine behind the normal form is coset enumeration along the chain
$S_1 < S_2 < \dots < S_n < S_{n+1}$, where $S_n$ sits inside $S_{n+1}$ as
the permutations fixing the bottom wire --- syntactically, the circuits
of the form $c\;{\uparrow}$.  Read as permutations, two circuits lie in the
same right coset of $S_n$ exactly when they bring the same wire to the
bottom position; there are $n{+}1$ choices, so $S_n$ has index $n{+}1$,
and a natural representative of each coset is the \emph{descending
staircase} of swaps that carries that wire down.  In the library a coset
is not a set of circuits but a label naming one, an element of the finite type
$\aid{C}\;n$; here the label is the length of the staircase, as a unary
numeral bounded by the width, and the \emph{section} $[\,\_\,]^{\mathsf
c}$ maps it to the staircase circuit.  The representatives on three,
two, and one wires are drawn on the right:

\begin{center}
\begin{minipage}[c]{0.30\linewidth}
\begingroup\setlength{\mathindent}{0pt}
\begin{code}
data C : ℕ → Set where
  ε   : C n
  σ•_ : C n → C (₁₊ n)
[_]ᶜ : C n → Circuit (₁₊ n)
[ ε ]ᶜ     = ε
[ σ• c ]ᶜ  = σ • ([ c ]ᶜ ↑)
\end{code}
\endgroup
\end{minipage}\hspace{2.5em}%
\begin{minipage}[c]{0.42\linewidth}\centering
\begin{tabular}{@{}l@{\;\;}ccc@{}}
  $\aid{C}\;2$: & \ufig[0.35]{c2-0} & \ufig[0.35]{c2-1} & \ufig[0.35]{c2-2}\\
  & $\varepsilon$ & $\sigma{\bullet}\:\varepsilon$ &
    $\sigma{\bullet}\:\sigma{\bullet}\:\varepsilon$\\[4pt]
  $\aid{C}\;1$, $\aid{C}\;0$: & \ufig[0.35]{c1-0} & \ufig[0.35]{c1-1} & \ufig[0.35]{c0-0}\\
  & $\varepsilon$ & $\sigma{\bullet}\:\varepsilon$ & $\varepsilon$
\end{tabular}
\end{minipage}
\end{center}

\noindent So $\aid{C}\;n$ has exactly the $n{+}1$ elements drawn above
for $n = 2, 1, 0$, matching the index of $S_n$ in $S_{n+1}$.  Why this yields
a normal form is the classical argument of computational group
theory~\cite{sims1994computation,holt2005handbook}:
every $f \in S_{n+1}$ lies in one coset, say with representative $c_n$,
so $f = f' \cdot c_n$ with $f' \in S_n$, and recursing on $f'$ writes $f$
uniquely as $c_1 c_2 \cdots c_n$ with $c_k \in \aid{C}\;k$ --- a numeral
in the \emph{factorial number system}, one digit per level, since
$|S_{n+1}| = (n{+}1) \cdot n \cdots 2$.

\emph{Computing} this factorisation is the job of the coset table.
Given a coset, represented by its staircase, and one generator, the
table says where their product lands --- the updated coset --- and what
is left over on the wires above --- the residual circuit --- so that
staircase followed by generator equals residual (shifted up) followed by
updated staircase.
Pushing $\sigma$ into a staircase can lengthen it, cancel with it ($\sigma^2 =
\varepsilon$), or pass it through and shift it up (Yang--Baxter), or
commute past; for a generator that does not touch the bottom wire we can
recurse.  The residual is empty or a single swap here, but in richer
gate sets it can be a longer circuit, which is why the table returns a
circuit on one wire fewer rather than a generator
(Figure~\ref{fig:pushing} draws the four cases):

\begin{center}
\begin{minipage}[c]{0.65\linewidth}
\begingroup\setlength{\mathindent}{0pt}
\begin{code}
ract : C (₁₊ n) → Gen (₂₊ n) → Circuit (₁₊ n) × C (₁₊ n)
ract ε σ-gen = ε , σ• ε             -- grows
ract (σ• ε) σ-gen = ε , ε           -- cancels (σ² = ε)
ract (σ• σ• c) σ-gen = σ , σ• σ• c  -- passes (Yang-Baxter)
ract ε (g ↥) = [ g ]ʷ , ε           -- commutes past
ract (σ• c) (g ↥) = proj₁ (ract c g) ↑ , σ• (proj₂ (ract c g))  -- recurses
\end{code}
\endgroup
\end{minipage}\hspace{0.5em}%
\begin{minipage}[c]{0.30\linewidth}\raggedright
  \ufig[0.35]{ract-l}\;\scalebox{0.7}{$\approx$}\;\ufig[0.35]{ract-r}
\end{minipage}
\end{center}

\begin{figure}[h]
\centering
\begin{tabular}{@{}r@{\,}c@{\,}l@{\;\;}l@{\qquad}r@{\,}c@{\,}l@{\;\;}l@{}}
  \ufig[0.35]{push4-l} & \scalebox{0.7}{$\rightarrow$} & \ufig[0.35]{push4-r} &
    \small grows &
  \ufig[0.35]{push3-l} & \scalebox{0.7}{$\rightarrow$} & \ufig[0.35]{push3-r} &
    \small cancels ($\sigma^2 = \varepsilon$)\\[4pt]
  \ufig[0.35]{push2-l} & \scalebox{0.7}{$\rightarrow$} & \ufig[0.35]{push2-r} &
    \small passes (Yang--Baxter) &
  \ufig[0.35]{push1-l} & \scalebox{0.7}{$\rightarrow$} & \ufig[0.35]{push1-r} &
    \small commutes past
\end{tabular}
\caption{Pushing a swap through a staircase: the four cases of the coset
table \aid{ract} as circuit rewrites.}
\label{fig:pushing}
\end{figure}

\medskip\noindent
\aid{σ-gen} abbreviates $\aid{gate₂}\;\aid{σ-gate}$.  The last clause is the recursion: a
shifted generator meets a staircase by meeting its tail one wire down.
Each clause is certified by the \emph{soundness law}, pictured on the
right and proved once per table by a short equational derivation from
the axioms:

\begin{code}
ract-sound : ∀ c b → let (b' , c') = ract c b in [ c ]ᶜ • [ b ]ʷ ≈ b' ↑ • [ c' ]ᶜ
\end{code}
\begin{code}[hide]
ract-sound = proof-elided
\end{code}

The stateful word traversal \aid{{\_}ᵗ} extends
\aid{ract} to whole circuits, threading the coset from left to right and
concatenating the residuals; run from the identity coset, it maps an
$(n{+}2)$-wire circuit to an $(n{+}1)$-wire circuit and a coset --- one
level of normalisation.

What the library checks a table against is fixed once and for all.
One level is a Reidemeister--Schreier situation: a subgroup
presentation $\Gamma$ and a group presentation $\Delta$, a finite type
$C$ of \emph{coset labels} with an identity coset $I$, a section
$[\_] : C \to \aid{Word}\;Y$ mapping each label to its representative,
the generator embedding $f$, and the coset table $h : C \to Y \to
\aid{Word}\;X \times C$, which pushes one letter of the group past a
coset representative and returns the residual word of the subgroup and
the updated coset.  The classical algorithms \emph{find} such a table;
the library takes a claimed table and \emph{verifies} it against five
hypotheses: (1) the table inverts the embedding at the identity coset;
(2) it respects $\Delta$'s axioms, coset by coset (\aid{h-wd-ax}); (3)
the embedding $f$ respects $\Gamma$'s axioms; (4) the identity coset is
represented by $\varepsilon$, $[\,I\,] \approx \varepsilon$; and (5)
the \emph{soundness law}, that sliding one letter past a representative
matches the table.  From these it derives the normal-form transport
$\aid{NormalForm}\;\Gamma\;\NF \to \aid{NormalForm}\;\Delta\;(\NF
\times C)$, and \aid{CosetTower} iterates it up an $\mathbb{N}$-indexed
family of presentations, one level per wire.  Here $\Gamma$ and
$\Delta$ are the relations on $n{+}1$ and $n{+}2$ wires, $f$ is the
shift, $h$ is \aid{ract}, and (5) is \aid{ract-sound}; (3) is the
structural rule \aid{cong↑} itself, (1) and (4) hold by evaluation, and
(2) is a case analysis over the axiom instances, coset by coset, whose
coset halves hold by evaluation and whose residual halves are one-line
derivations (the unit laws, or the axiom being pushed).  For the
symplectic group even the coset half of (2) does not close by
evaluation, and \S\ref{sec:pushing} explains what replaces it.

\emph{Iterating} the level is the coset tower, whose carrier grows by
one digit per level:

\begin{center}
\begin{minipage}[c]{0.56\linewidth}
\begingroup\setlength{\mathindent}{0pt}
\begin{code}
NF : ℕ → Set
NF 0       = ⊤
NF (suc k) = NF k × C k   -- ⊤ × C 0 × C 1 × ⋯ × C k
\end{code}
\endgroup
\end{minipage}\hspace{2.5em}%
\begin{minipage}[c]{0.16\linewidth}\centering
  \ufig[0.35]{fig3_staircase}
\end{minipage}
\end{center}

\noindent so $\aid{NF}\;n$ has exactly $n!$ elements.  The normal-form
function $\aid{nf} : \aid{Circuit}\;n \to \aid{NF}\;n$ runs the tables;
its \emph{section} $\aid{inv-nf} : \aid{NF}\;n \to \aid{Circuit}\;n$
rebuilds a circuit from its digits by concatenating the staircases,
suitably shifted, as drawn on the right for $(\mathsf{tt}, c_0, c_1,
c_2, c_3) \in \aid{NF}\;4$, the digit $c_0 \in \aid{C}\;0$ being the
empty staircase --- one staircase per level, each shifted up to the
wires of its level: a ``staircase of staircases''.

\subsection{Semantics, uniqueness, and completeness}\label{sec:sn-semantics}

So far everything was syntax.  The development interprets circuits by a
semantics $\aid{⟦\_⟧} : \aid{Circuit}\;n \to \aid{Permutation′}\;n$
into the standard library's permutations of $\mathsf{Fin}\;n$
($\sigma$ swaps the bottom two elements), bundled as a \aid{Group},
$\aid{Permutation′-group}\;n$, by our \texttt{ForStdlib} layer.
\emph{Soundness} --- $w \approx v$
implies $\aid{⟦}w\aid{⟧} \doteq \aid{⟦}v\aid{⟧}$ --- is a small case
analysis over the axioms.  The heart of the matter is the \emph{unique
normal form} property, that normal forms with
equal denotations are equal: $\aid{⟦}\;\aid{inv-nf}\;u\,\aid{⟧} \doteq
\aid{⟦}\;\aid{inv-nf}\;v\,\aid{⟧}$ implies $u \equiv v$, proved by
induction on the tower: the lower digits fix the last point, so the top
digit alone determines where the permutation sends it; equal denotations
therefore force equal top digits, and cancelling them reduces the claim
to the level below.  Soundness and uniqueness are exactly the premises of the generic
lemma \aid{by-normalization}, which concludes \emph{completeness}:
$\aid{⟦}w\aid{⟧} \doteq \aid{⟦}v\aid{⟧}$ implies $w \approx v$.
Unfolding the lemma on this example shows where the axioms enter.  Each
step of \aid{nf} is one of the rewrites of Figure~\ref{fig:pushing}, and
each of those is an instance of \aid{order}, of \aid{yang-baxter}, or of
the structural rules (the far commutation \aid{comm₂}, lifted to the
wires above by \aid{cong↑}); so $w \approx \aid{inv-nf}\;(\aid{nf}\;w)$
--- the law of the \aid{NormalForm} witness --- is a derivation from
exactly those rules.  Now suppose $\aid{⟦}w\aid{⟧} \doteq
\aid{⟦}v\aid{⟧}$.  By soundness the normal forms of $w$ and $v$ have the
same denotation, so by uniqueness they are equal, and hence $w \approx
\aid{inv-nf}\;(\aid{nf}\;w) = \aid{inv-nf}\;(\aid{nf}\;v) \approx v$, a
derivation assembled from the two normal-form derivations and nothing
else.  Completeness thus never appeals to a rule beyond the axioms and
structural rules above.  \texttt{Examples.Groups.Symmetric} states
these for every $n$.  Soundness, completeness, and surjectivity (every
permutation is a product of adjacent transpositions) assemble into the
library's \emph{presentation record}.
The record packages two things: a \aid{Grouplike} witness, which makes
the presented monoid a group, and a proof that the semantics
$\aid{⟦\_⟧}$ is a group isomorphism from that group onto the target.  We quote
this one statement in Agda because every presentation theorem in the
paper has its shape:

\begin{code}
symmetric-presentation : ∀ n → (n VRel,_===_) IsPresentationOf (Permutation′-group n)
\end{code}
\begin{code}[hide]
symmetric-presentation n = SymPres.presentation {n}
\end{code}

\noindent read: the monoid of swap circuits modulo \{\aid{order},
\aid{yang-baxter}\} and the structural rules \emph{is} the group of
permutations of $\mathsf{Fin}\;n$ (\aid{IsPresentationOf} is the mixfix
record name \aid{{\_}IsPresentationOf{\_}} used infix).  The symplectic
factor of \S\ref{sec:symplectic} is this pipeline again, with richer
gates and a richer action, and \S\ref{sec:composing} composes such
presentations.
